\section{Limitations and future work}
\label{sec:limitations}

The fragment $\CICm$ was chosen to make every emitted axiom provably
sound while covering the phenomena the factored architecture is about:
pattern matching, structural recursion, propositional content in types,
singleton eliminations.  We now delimit the boundary, distinguishing
presentational restrictions (liftable with more bookkeeping) from
genuine semantic boundaries.

\subsection{Polymorphism and monomorphisation}
\label{sec:limitations-poly}

$\CICm$ has no type variables: every inductive type is declared at a
ground instance, and guards are unary per-datatype predicates.  This
matches the intended pipeline, in which heuristic monomorphisation ---
instantiating polymorphic definitions and inductive families at the
ground instances relevant to a problem, as HOL hammers do
\cite{BlanchetteEtAl2016} --- runs before the translation of this paper.
Extending the soundness proof to genuinely polymorphic environments
(type-variable guards, parameterised guard predicates) is orthogonal
future work; the design questions there (which type arguments and guards
are omissible) are treated in companion notes to this paper and in the
monotonicity literature \cite{BlanchetteEtAl2016}.

\subsection{Indexed families and large elimination}
\label{sec:limitations-indexed}

Datatypes in $\CICm$ are non-indexed (inductive \emph{predicates} carry
indices, but their proofs are erased).  Indexed informative families
(vectors, finite sets) require dependent case motives over indices and
transport in branches; the semantic machinery (classes, stages, ranks)
extends, but erasure of index information interacts with injectivity in
ways that need care, and we have not carried it out.  Large elimination
(types defined by recursion) is excluded outright: it breaks the rigidity
of set types (\cref{lem:rigid}) on which the structural definitions of
$\guardf$ and $\semr{\cdot}{}$ rest.  A translation for large-eliminating
code would need a different device (e.g.\ per-instance specialisation).

\subsection{Well-founded recursion: a hard constraint}
\label{sec:limitations-wf}

$\tfix$ captures structural recursion only.  The extension to
well-founded ($\mathsf{Acc}$-based) recursion is \emph{not} a routine
generalisation, and the obstruction deserves precise statement, since it
constrains any implementation of the architecture.

In CIC, a function defined by well-founded recursion --- Euclidean
division $\mathsf{div} : \forall a\,b.\ b \neq 0 \to \{q \mid \dots\}$
is the canonical example --- recurses through a proof of
$\mathsf{Acc}(x)$, and its unfolding equation holds \emph{only under the
erased propositional preconditions}.  Program extraction may erase the
precondition entirely: the extracted $\mathsf{div}$ diverges at $b = 0$,
which is harmless because lazy (or guarded) evaluation never forces the
offending closure \cite[\S 2.6]{Letouzey2004,Letouzey2003}.  An
equational axiomatisation enjoys no such protection: an equation is
``already applied everywhere''.  Erasing the precondition from the
unfolding equation of $\mathsf{div}$ yields the instance
\[
  \Ihat{\mathsf{div}}(a, \csO) \feq \csS(\Ihat{\mathsf{div}}(a, \csO))
  \quad\text{(for } b = \csO\text{, when the guard test } b \leq a
  \text{ degenerates)},
\]
which contradicts the perfectly legitimate premise $\forall n.\ n \neq
\csS(n)$ --- the resulting theory is \emph{inconsistent}, not merely
incomplete, and \cref{cor:consistency} would fail.  The constraint is
therefore: \emph{defining equations obtained by recursion through an
erased proof must retain the erased propositional premises of the
definition as implications} ($\forall a\,b.\ \F{b \neq \csO} \to
\dots$), the analogue for equations of the evaluation-order problem
that forced Letouzey to abandon the erasure rule
$\erase{\lambda p{:}\varphi.t} = \erase{t}$ for \emph{executable}
extraction.  In our fragment the unconditional-drop rule of
\cref{def:erasure} is sound precisely because $\tfix$ types contain no
propositional premises and unfolding is total on the guard datatype
(\cref{lem:fundamental}); lifting either restriction requires the
premised-equation form, with the semantics interpreting the defined
symbol by an arbitrary total extension of the partial function.

\subsection{Transport erasure and UIP}
\label{sec:limitations-transport}

The clause $\erase{\tcast{\pi}{z.\tau}{t}} = \erase{t}$ (and the
corresponding validity of unconditional transport equations, e.g.\ the
image of $\mathsf{eq\_rect}$ becoming an identity) is validated by $\ML$
through proof irrelevance: $\pv{a = b} = 1$ collapses the classes, so
transport is semantically the identity.  As a \emph{source} theorem,
however, ``transport is the identity'' at an arbitrary type is
uniqueness of identity proofs, which CIC does not prove
\cite{GilbertEtAl2019}.  The consequence is not unsoundness --- the
model validates it, and $\CICm$'s $\tcast{}{}{}$ is well typed only at
matching endpoints --- but a \emph{reconstruction} caveat for the
implemented pipeline: an ATP proof essentially using transport-erasure
at a type without decidable equality may resist replay inside Coq, the
same failure mode the tool already accepts for other
semantically-validated principles.  A conservative variant guards the
transport equation with the source equation as a premise.

\subsection{Function extensionality fails in the term model}
\label{sec:limitations-funext}

$\ML$ interprets equality as $\lbox$-convertibility of closed terms,
which is intensional: $\cls{\lambda x.\,\mathsf{add}(\csO, x)} \neq
\cls{\lambda x.\,x}$ although the two agree on all of $\semr{\nat}{}$
(the first is stuck at a variable).  So $\pv{\text{funext}} = 0$: the
guarantee of \cref{thm:soundness} does not extend to environments
postulating function extensionality (\cref{rem:axioms}).  Repairing
this requires a coarser model --- quotienting $\Dom$ by an
observational equivalence --- which changes the discrimination and
compilation arguments and is left as future work.  Note the asymmetry:
the translation never \emph{emits} an extensionality axiom, so
\cref{cor:consistency} is unaffected; only premise sets containing
funext-derived lemmas step outside the present soundness statement.

\subsection{Other presentational restrictions}
\label{sec:limitations-misc}

Mutual and nested inductive types, deep (nested-pattern) and mutual
recursion, and fixpoints whose recursive calls occur under inner binders
with non-variable guard arguments are excluded by the syntactic shape of
$\tfix$ and the declarations; all appear liftable by elaborating them
into the present forms (mutual recursion via a sum type, deep guards via
nested case trees --- which the \emph{compilation} already handles,
cf.\ \cref{ex:compile}(3) --- and course-of-values recursion via
auxiliary structural functions), at bookkeeping cost in the fix case of
\cref{lem:fundamental}.  Inductive predicates with predicate
\emph{parameters} (e.g.\ $\mathsf{Forall}$) are excluded from $\Sfrag$;
per-instance specialisation (a fresh predicate symbol per concrete
predicate parameter, with clauses inlined) extends the translation
shallowly and is the natural implementation route.  Enumeration-like
datatypes whose constructors carry only propositional data
($\mathsf{sumbool}$) are covered as plain datatypes here; inlining
their inversion disjunction into guard positions is a further
shallowness optimisation compatible with the model.

\subsection{Completeness}
\label{sec:limitations-completeness}

The translation is not complete, deliberately.  Beyond the classical
semantics (\cref{rem:not-claimed}), first-order provability from $\ThE$
lacks: induction (only case analysis is emitted; goals needing
induction are the reconstruction backend's or the user's job),
conversion under binders for lifted symbols (two extensionally equal
lifted functions get unrelated symbols --- the $\xi$-rule failure
discussed in \cite[Example 51]{Czajka2018}), and anything involving
non-shallow propositions.  A precise completeness theorem for a
quantifier fragment (e.g.\ universal equational goals over first-order
datatypes, where \cref{prop:standard} makes $\ML$ standard) would be
valuable; we expect the per-constructor equations to be complete for
ground equational consequences of the definitions, by
\cref{rem:eq-shape} and standard rewriting arguments
\cite{BaaderNipkow1998}, but leave the statement to future work.

\subsection{Proof-theoretic soundness and implementation}
\label{sec:limitations-proof-theoretic}

Two further directions.  First, a proof-theoretic soundness theorem in
the style of \cite{Czajka2018} --- transforming a first-order
derivation from $\ThE$ into a $\CICm$ proof term, presumably modulo
excluded middle and proof irrelevance --- would sharpen
\cref{thm:soundness} into a reconstruction algorithm; the
substitution-stability machinery needed for lifted symbols is exactly
the hard part there, and our semantic treatment sidesteps rather than
solves it.  Second, on the implementation side, the axiom families of
\cref{def:theory} map one-to-one onto sites in the existing CoqHammer
translation (case lifting, guard generation, definition axioms), so the
architecture can be adopted incrementally --- per-constructor equations
first, then singleton erasure and refinement expansion, then the full
specification extraction --- with each stage independently measurable
against the current translation on the standard evaluation harness.
